\documentclass[11pt]{amsart}

\usepackage[T1]{fontenc}
\usepackage{lmodern}
\usepackage{microtype}
\usepackage{amsmath,amssymb,amsthm}
\usepackage[hidelinks]{hyperref}

\newtheorem{theorem}{Theorem}
\newtheorem{proposition}[theorem]{Proposition}
\newtheorem{lemma}[theorem]{Lemma}
\theoremstyle{remark}
\newtheorem{remark}[theorem]{Remark}

\DeclareMathOperator{\ill}{ill}
\DeclareMathOperator{\ext}{ext}
\DeclareMathOperator{\Int}{Int}
\DeclareMathOperator{\conv}{conv}
\newcommand{\CC}{\mathbb{C}}
\newcommand{\RR}{\mathbb{R}}
\newcommand{\illstar}{\ill^{*}}

\title[The complex cross-polytope at $n=2$]{The Complex Cross-Polytope
\texorpdfstring{$B_{1}(\CC^{2})$}{B1(C2)}: Real Fractional Illumination Number
\texorpdfstring{$2^{2}$}{2^2}, and Not a Complex-Linear Image of the Polydisc}
\date{}
\subjclass[2020]{Primary 52C17; Secondary 52A21, 52A37, 52A20}
\keywords{complex convex body, illumination, fractional illumination number,
weighted covering, polydisc, complex zonoid, cross-polytope, counterexample}

\hypersetup{
  pdftitle={The complex cross-polytope B_1(C^2): real fractional illumination
            number 2^2, and not a complex-linear image of the polydisc}
}

\begin{document}

\begin{abstract}
Rotem, Schejter and Slomka conjectured that every complex convex body
$K\subseteq\CC^{n}$ satisfies $\illstar(K)\le 2^{n}$, with equality if and only
if $K$ is a linear image of the polydisc $D^{n}$. Reading $\illstar$ as the
ordinary real fractional illumination number of $K$ regarded as a convex body in
$\RR^{2n}$ --- a reading we do not check against the conjecture's own wording of
$\illstar$ --- the complex cross-polytope
$B_{1}(\CC^{2})=\{z:|z_{1}|+|z_{2}|\le 1\}$ satisfies
$\illstar(B_{1}(\CC^{2}))=4=2^{2}$, and it is not a complex-linear image of
$D^{2}$. So under that reading of $\illstar$, and with ``linear image'' read as
complex-linear, the equality clause fails at $n=2$. The bound
$\illstar(K)\le 2^{n}$ itself is not contradicted --- the witness attains
it with equality --- and the equality clause is untouched for $n\ge 3$, because
$\illstar(B_{1}(\CC^{n}))=2n<2^{n}$ there.
\end{abstract}

\maketitle

\section{The conjecture}\label{sec:conj}

Let $D\subseteq\CC$ be the closed unit disc and $D^{n}=D\times\dots\times D$ the
polydisc. Following Rotem, Schejter and Slomka \cite{RSS}, a \emph{complex
convex body} in $\CC^{n}$ is the unit ball of a norm on $\CC^{n}$; equivalently
it is a centrally symmetric convex body $K\subseteq\RR^{2n}$ with
$x\in K\Rightarrow e^{i\theta}x\in K$ for all $\theta\in\RR$. A direction
$v$ \emph{illuminates} a boundary point $x$ of a convex body $K$ if
$x+tv\in\Int K$ for some $t>0$; writing $A_{K}(x)$ for the set of directions
illuminating $x$, a non-negative Borel measure $\mu$ \emph{illuminates} $K$ if
$\mu(A_{K}(x))\ge 1$ for every $x\in K$, and
\[
  \illstar(K)=\inf\bigl\{\mu(\RR^{2n})\;:\;\mu\text{ illuminates }K\bigr\}
\]
is the \emph{fractional illumination number} of $K$, introduced by
Nasz\'odi \cite{Naszodi}. Throughout, $\illstar$ of a complex body is read as the
ordinary real fractional illumination number of $K$ regarded as a convex body in
$\RR^{2n}$: real directions, and no circle invariance imposed on $\mu$. We do not
reproduce \cite{RSS}'s own wording of that definition here, and every value below
is asserted for this reading only; under a reading that restricted $\mu$ to
circle-invariant measures the numbers could differ. The measure we exhibit below
is supported on unit vectors, so it is admissible whether or not non-unit vectors
are allowed. Everything below is therefore conditional on \cite{RSS}'s $\illstar$
agreeing with the reading just fixed; we have not established that it does.

The statement we address is Conjecture 1.2 of \cite{RSS}. In the version of
record it reads, verbatim:

\begin{quote}
\textbf{Conjecture 1.2.} Given a complex convex body $K\subseteq\CC^{n}$, we
have $\ill^{*}(K)\le 2^{n}$. Moreover, equality holds if and only if $K$ is a
linear image of $D^{n}$.
\end{quote}

\noindent
We quote the published wording; the companion Conjecture 1.1 of \cite{RSS}
concerns the classical illumination number $\ill$ and is \emph{not} addressed
here.

Two theorems of \cite{RSS} delimit what was known: $\illstar(D^{n})=2^{n}$ for
all $n\ge1$, and every full-dimensional complex zonoid
$Z\subseteq\CC^{n}$ that is not a linear image of $D^{n}$ has
$\illstar(Z)<\illstar(D^{n})$. By the second of those theorems, any
counterexample to the equality clause must lie outside the complex zonoids;
Section~\ref{sec:zonoid} shows that our witness, the polar of the polydisc,
does.

\section{The result}\label{sec:result}

\begin{theorem}\label{thm:main}
Let $K=B_{1}(\CC^{2})=\{z\in\CC^{2}:|z_{1}|+|z_{2}|\le1\}$. With $\illstar$ read
as in Section~\ref{sec:conj},
\[
  \illstar\bigl(B_{1}(\CC^{2})\bigr)=4=2^{2},
\]
and $B_{1}(\CC^{2})$ is not a complex-linear image of $D^{2}$. Consequently, if
$\illstar$ in Conjecture 1.2 of \textup{\cite{RSS}} is that same number and its
``linear image'' is meant complex-linearly, the equality clause fails at $n=2$:
the bodies in $\CC^{2}$ attaining $\illstar=2^{2}$ include at least the two
$GL_{2}(\CC)$-orbits of $D^{2}$ and of $B_{1}(\CC^{2})$.
\end{theorem}

We state at once what is \emph{not} claimed, since the conjecture is a
conjunction and only one conjunct falls.

\begin{itemize}
\item The two readings fixed in Section~\ref{sec:conj} are not shown to be
      \cite{RSS}'s. We have not reproduced their definition of $\illstar$, and we
      settle only complex-linear images of $D^{2}$; whether $B_{1}(\CC^{2})$ is a
      real-linear image of $D^{2}$ is not decided here.
\item The bound $\illstar(K)\le2^{n}$ is \emph{not} contradicted. The witness
      satisfies it with equality, and we have no evidence against it; the first
      conjunct of Conjecture 1.2 remains open.
\item Nothing is claimed for $n\ge3$. By Proposition~\ref{prop:dim} below,
      $\illstar(B_{1}(\CC^{n}))=2n<2^{n}$ for $n\ge3$, so this witness is silent
      there and the equality clause may well hold as written for every $n\ge3$.
\item Even at $n=2$ we show only that the equality set is strictly larger than
      the orbit of $D^{2}$. Whether $\{D^{2},B_{1}(\CC^{2})\}$ exhausts it is
      open.
\item Conjecture 1.1 of \cite{RSS}, on the classical illumination number, is
      untouched; we make no claim about it.
\end{itemize}

\section{The witness}\label{sec:witness}

Throughout, $K=B_{1}(\CC^{2})$, and we use the real pairing
$\langle u,z\rangle=\operatorname{Re}(u_{1}\overline{z_{1}}+u_{2}\overline{z_{2}})$
on $\CC^{2}\cong\RR^{4}$.

\smallskip
\noindent\textbf{(a) $K$ is a complex convex body.} $\|z\|:=|z_{1}|+|z_{2}|$ is
a norm on $\CC^{2}$ and $K$ is its unit ball; $K$ is full-dimensional (real
dimension $4$) and $\|e^{i\theta}z\|=\|z\|$, so $K$ is circle-invariant.
Moreover
\[
  K=\conv\bigl(De_{1}\cup De_{2}\bigr)=(D^{2})^{\circ},
\]
the first equality because $z=|z_{1}|\bigl(\tfrac{z_{1}}{|z_{1}|},0\bigr)
+|z_{2}|\bigl(0,\tfrac{z_{2}}{|z_{2}|}\bigr)+(1-\|z\|)\cdot 0$ is a convex
combination for every $z\in K$.

\smallskip
\noindent\textbf{(b) Support function.}
\[
  h_{K}(u)=\sup_{|z_{1}|+|z_{2}|\le1}\langle u,z\rangle
         =\sup_{0\le s\le1}\bigl(|u_{1}|s+|u_{2}|(1-s)\bigr)
         =\max(|u_{1}|,|u_{2}|),
\]
the complex $\ell^{\infty}$ norm --- the polar situation to $D^{2}$, whose
support function is $|u_{1}|+|u_{2}|$.

\smallskip
\noindent\textbf{(c) The boundary, exhaustively.} $\partial K$ consists of
exactly three classes: (i) $|z_{1}|=1,\ z_{2}=0$; (ii) $z_{1}=0,\ |z_{2}|=1$;
(iii) \emph{mixed} points, with both coordinates nonzero and
$|z_{1}|+|z_{2}|=1$.

\smallskip
\noindent\textbf{(d) Extreme points: two disjoint circles.}
\[
  \ext(K)=C_{1}\cup C_{2},\qquad
  C_{1}=\{(e^{i\alpha},0)\},\qquad C_{2}=\{(0,e^{i\beta})\}.
\]
Indeed $\ext(\conv(A\cup B))\subseteq\ext(A)\cup\ext(B)$ with $A=De_{1}$,
$B=De_{2}$ gives the inclusion, and each such point is exposed: on $K$,
$\operatorname{Re}(e^{-i\alpha}z_{1})\le|z_{1}|\le|z_{1}|+|z_{2}|\le1$ with
equality if and only if $z_{1}=e^{i\alpha}$ and $z_{2}=0$.

\smallskip
\noindent\textbf{(e) The illuminating measure, in closed form.}
\begin{equation}\label{eq:mu}
  \mu \;=\; 2\cdot\operatorname{Unif}\{(e^{i\varphi},0):\varphi\in[0,2\pi)\}
        \;+\; 2\cdot\operatorname{Unif}\{(0,e^{i\psi}):\psi\in[0,2\pi)\},
  \qquad \mu(\RR^{4})=4 .
\end{equation}
It is supported on unit vectors.

\smallskip
\noindent\textbf{(f) Illumination criteria.} Since
$\|(e^{i\alpha},0)+tv\|=|e^{i\alpha}+tv_{1}|+t|v_{2}|
 =1+t\bigl(\operatorname{Re}(e^{-i\alpha}v_{1})+|v_{2}|\bigr)+O(t^{2})$,
\begin{equation}\label{eq:crit}
\begin{aligned}
  v\text{ illuminates }(e^{i\alpha},0)
    &\iff \operatorname{Re}(e^{-i\alpha}v_{1})+|v_{2}|<0,\\
  v\text{ illuminates }(0,e^{i\beta})
    &\iff \operatorname{Re}(e^{-i\beta}v_{2})+|v_{1}|<0,\\
  v\text{ illuminates a mixed }x
    &\iff \operatorname{Re}(e^{-i\arg x_{1}}v_{1})
          +\operatorname{Re}(e^{-i\arg x_{2}}v_{2})<0 .
\end{aligned}
\end{equation}
The same criteria follow from the normal-cone description used in \cite{RSS}
(namely, $v$ illuminates $x$ iff $\langle u,v\rangle<0$ for every nonzero outer
normal $u$ at $x$, cf.\ \cite[(2.2)]{Schneider}): at $x=(e^{i\alpha},0)$ the
normal cone is $\{(re^{i\alpha},u_{2}):r\ge|u_{2}|\}$, and at a mixed $x$ it is
the single ray spanned by $(e^{i\arg x_{1}},e^{i\arg x_{2}})$.

\section{Proof that \texorpdfstring{$\illstar(B_{1}(\CC^{2}))=4$}{ill*=4}}
\label{sec:value}

\begin{proof}[Lower bound $\illstar(K)\ge4$]
\emph{Step 1: the two extreme circles are illumination-disjoint.} If $v$
illuminates some point of $C_{1}$ then, by \eqref{eq:crit},
$\operatorname{Re}(e^{-i\alpha}v_{1})<-|v_{2}|$, whence
$|v_{1}|\ge-\operatorname{Re}(e^{-i\alpha}v_{1})>|v_{2}|$. Symmetrically,
illuminating a point of $C_{2}$ forces $|v_{2}|>|v_{1}|$. Hence
$A_{K}(C_{1})\subseteq V_{1}:=\{|v_{1}|>|v_{2}|\}$ and
$A_{K}(C_{2})\subseteq V_{2}:=\{|v_{2}|>|v_{1}|\}$, and $V_{1}\cap V_{2}
=\emptyset$: no single direction serves both circles.

\emph{Step 2: exact arc length.} Fix $v\in V_{1}$ and write $v_{1}=re^{i\varphi}$
with $r>0$ and $c=|v_{2}|/r\in[0,1)$. By \eqref{eq:crit}, $v$ illuminates
$x_{\alpha}:=(e^{i\alpha},0)$ if and only if $\cos(\varphi-\alpha)<-c$, an arc of
$\alpha$ whose endpoints satisfy $\cos(\varphi-\alpha)=-c$ and whose length is
exactly
\[
  2\pi-2\arccos(-c)=2\arccos(c)\;\le\;\pi ,
\]
with length $\pi$ only when $v_{2}=0$. (If $r=0$ the condition reads
$|v_{2}|<0$, which never holds.)

\emph{Step 3: averaging.} Let $\mu$ illuminate $K$. Integrating the constraint
$\mu(A_{K}(x_{\alpha}))\ge1$ over $\alpha\in[0,2\pi)$ and exchanging the order of
integration (Tonelli; all integrands are non-negative),
\[
  2\pi\;\le\;\int_{0}^{2\pi}\mu\bigl(A_{K}(x_{\alpha})\bigr)\,d\alpha
       \;=\;\int_{V_{1}}2\arccos\!\bigl(|v_{2}|/|v_{1}|\bigr)\,d\mu
       \;\le\;\pi\,\mu(V_{1}),
\]
so $\mu(V_{1})\ge2$, and identically $\mu(V_{2})\ge2$. The two sets are
disjoint, so $\mu(\RR^{4})\ge4$.
\end{proof}

\begin{proof}[Upper bound $\illstar(K)\le4$]
Take $\mu$ as in \eqref{eq:mu} and check the constraint on every boundary class
directly, so that no extreme-point reduction is needed.
\begin{itemize}
\item $x=(e^{i\alpha},0)$: by \eqref{eq:crit}, $(e^{i\varphi},0)$ illuminates $x$
      iff $\cos(\varphi-\alpha)<0$, an open half-circle, of measure exactly $1$
      under the mass-$2$ uniform measure on that circle; and $(0,e^{i\psi})$ never
      does, since the criterion would read $|e^{i\psi}|=1<0$. Total
      $\mu(A_{K}(x))=1$.
\item $x=(0,e^{i\beta})$: symmetric, total $1$.
\item $x$ mixed: each axis circle contributes its open half-circle, total
      $1+1=2$.
\item $x\in\Int K$: every direction illuminates $x$, total $4$.
\end{itemize}
Hence $\mu$ illuminates $K$ and $\illstar(K)\le\mu(\RR^{4})=4$.
\end{proof}

Combining the two bounds gives $\illstar(B_{1}(\CC^{2}))=4=2^{2}$.

\section{\texorpdfstring{$B_{1}(\CC^{2})$}{B1(C2)} is not a complex zonoid, hence
not a complex-linear image of \texorpdfstring{$D^{2}$}{D2}}\label{sec:zonoid}

A complex-linear image of $D^{2}$ is a $2$-generator complex zonotope, since
$T(D^{2})=T(De_{1}+De_{2})=D\,T(e_{1})+D\,T(e_{2})$; in particular it is a
complex zonoid. So Theorem~\ref{thm:main}'s second assertion follows from:

\begin{proposition}\label{prop:zonoid}
$B:=B_{1}(\CC^{2})$ is not a complex zonoid.
\end{proposition}

\begin{proof}[Proof: the mean-value obstruction]
Recall from \cite{RSS} that a complex zonoid $Z\subseteq\CC^{n}$ admits a
non-negative Borel measure $\nu$ on the unit sphere $S$ of $\CC^{n}$ with
$h_{Z}(\theta)=\int_{S}|\langle x,\theta\rangle_{\CC}|\,d\nu(x)$. Suppose such a
$\nu$ existed for $B$, whose support function is
$h_{B}(\theta)=\max(|\theta_{1}|,|\theta_{2}|)$. Testing $\theta=(1,0)$ and
$\theta=(0,1)$ gives
\begin{equation}\label{eq:12}
  \int_{S}|x_{1}|\,d\nu=1,\qquad \int_{S}|x_{2}|\,d\nu=1 ,
\end{equation}
and testing $\theta=(1,e^{it})$ for every real $t$ gives
$\int_{S}|x_{1}+e^{it}x_{2}|\,d\nu=1$, since $h_{B}=\max(1,1)=1$. Averaging the
last identity over $t\in[0,2\pi)$ and exchanging the order of integration
(Tonelli again),
\[
  \int_{S}M\bigl(|x_{1}|,|x_{2}|\bigr)\,d\nu=1,\qquad
  M(a,b):=\frac{1}{2\pi}\int_{0}^{2\pi}\bigl|a+e^{it}b\bigr|\,dt .
\]
The map $w\mapsto a+wb$ is holomorphic on the unit disc, so
$\frac{1}{2\pi}\int_{0}^{2\pi}(a+e^{it}b)\,dt=a$ and hence $M(a,b)\ge|a|$;
writing $a+e^{it}b=e^{it}(b+e^{-it}a)$ gives $M(a,b)\ge|b|$ as well. Therefore
$M(a,b)\ge\max(a,b)\ge(a+b)/2$ for $a,b\ge0$, and with \eqref{eq:12}
\[
  1=\int_{S}M\,d\nu\;\ge\;\int_{S}\frac{|x_{1}|+|x_{2}|}{2}\,d\nu
   \;=\;\frac{1+1}{2}\;=\;1 .
\]
Equality is forced, so $M(|x_{1}|,|x_{2}|)=(|x_{1}|+|x_{2}|)/2$ for
$\nu$-almost every $x$. But this is impossible on $S$, and the following three
cases are exhaustive there:
\begin{itemize}
\item if $|x_{1}|,|x_{2}|>0$ with $|x_{1}|\ne|x_{2}|$, no property of $M$ beyond
      $M\ge\max$ is needed: the second inequality of \,$M\ge\max(a,b)\ge(a+b)/2$
      is already strict, since $\max(a,b)>(a+b)/2$ whenever $a\ne b$;
\item if $|x_{1}|=|x_{2}|=r>0$ then $\max(a,b)=(a+b)/2$, so this is the one case
      in which the \emph{value} of $M$ must decide it, and it does:
      $M=r\cdot\frac{1}{2\pi}\int_{0}^{2\pi}2|\cos(t/2)|\,dt=\frac{4r}{\pi}
      =1.2732\ldots r>r=\max(r,r)$, using
      $|a+e^{it}a|=2a|\cos(t/2)|$ and $\int_{0}^{2\pi}|\cos(t/2)|\,dt=4$;
\item if $x_{2}=0$ then $M=|x_{1}|$ while $(|x_{1}|+0)/2=|x_{1}|/2$, forcing
      $x_{1}=0$, impossible on $S$; symmetrically if $x_{1}=0$.
\end{itemize}
In every case $M(|x_{1}|,|x_{2}|)>(|x_{1}|+|x_{2}|)/2$ strictly, so $\nu$ charges
no point of $S$.
Hence $\nu=0$, contradicting \eqref{eq:12}.
\end{proof}

\begin{remark}\label{rem:consequences}
No theorem of \cite{RSS} is contradicted: the theorem that a
full-dimensional complex zonoid other than a linear image of $D^{n}$ has
$\illstar$ \emph{strictly} below $\illstar(D^{n})$ does not apply to our
witness, precisely because of Proposition~\ref{prop:zonoid}.
\end{remark}

\section{The witness in \texorpdfstring{$\CC^{n}$}{Cn}}\label{sec:dim}

\begin{proposition}\label{prop:dim}
$\illstar\bigl(B_{1}(\CC^{n})\bigr)=2n$ for every $n\ge1$.
\end{proposition}

\begin{proof}
As in \eqref{eq:crit}, $v$ illuminates $e^{i\alpha}e_{i}$ if and only if
$\operatorname{Re}(e^{-i\alpha}v_{i})+\sum_{j\ne i}|v_{j}|<0$, which forces
$|v_{i}|>\sum_{j\ne i}|v_{j}|$; adding two such inequalities for $i\ne k$ yields
a contradiction, so the $n$ extreme circles are pairwise illumination-disjoint.
The argument of \S\ref{sec:value} then gives mass at least $2$ per circle, and
mass exactly $2$ per circle suffices: a boundary point whose set of nonzero
coordinates is $S$ receives exactly $|S|\ge1$.
\end{proof}

Now $2n=2^{n}$ holds only at $n=1$ and $n=2$, and $2n<2^{n}$ for $n\ge3$. So
this witness attains the conjectured extremal value only in $\CC^{2}$ and lies
strictly below it for $n\ge3$, where the equality clause is untouched by it. We
have not carried out a systematic priority check for the formula
$\illstar(B_{1}(\CC^{n}))=2n$ and make no claim that it is new; it may well be
recorded elsewhere.

\section{The accompanying program}\label{sec:program}

All proofs above are closed form. The program \texttt{verify.py} in this folder,
with its recorded run \texttt{verify.output.txt}, is a sampled consistency check
in exact rational arithmetic: it parses the body, the support function and the
measure \eqref{eq:mu}, tests the criteria \eqref{eq:crit} against the definition
of illumination (membership of $x+tv$ in $\Int K$) on finite samples of rational
unit vectors, and confirms $\mu(A_{K}(x))\ge1$ at sampled points of each boundary
class. It establishes none of the universal statements on which the theorem rests
--- the criteria at \emph{every} boundary point, $\ext(K)=C_{1}\cup C_{2}$, the
constraint $\mu(A_{K}(x))\ge1$ for \emph{every} $x$, and
$\illstar(B_{1}(\CC^{n}))=2n$ --- it hard-codes the cone mass $\mu(V_{1})\ge2$ of the lower bound rather than deriving
it, and its second route to the non-zonoid conclusion imports \cite{RSS}'s
zonoid face lemma. Those steps are the closed-form arguments printed above; the
program is not an independent verification of the theorem.

\begin{thebibliography}{9}

\bibitem{RSS}
L.~Rotem, B.~Schejter and B.~A.~Slomka,
\emph{The complex illumination problem},
Combinatorica \textbf{46} (2026), no.~1, Paper No.~3.
\href{https://doi.org/10.1007/s00493-025-00195-7}{doi:10.1007/s00493-025-00195-7};
e-print \href{https://arxiv.org/abs/2410.12021}{arXiv:2410.12021}.

\bibitem{Naszodi}
M.~Nasz\'odi,
\emph{Fractional illumination of convex bodies},
Contrib. Discrete Math. \textbf{4} (2009), no.~2, 83--88.

\bibitem{Schneider}
R.~Schneider,
\emph{Convex Bodies: The Brunn--Minkowski Theory},
2nd ed., Encyclopedia of Mathematics and its Applications \textbf{151},
Cambridge University Press, 2013.

\end{thebibliography}

\end{document}
